\documentclass[exercice,a4]{thedocument}%
\usepackage[french]{babel}%
\usepackage{pstricks}%
\usepackage{pst-plot}%
\usepackage{enumitem}%
\startdatakeys
    \source{21GENMATAA2}
    \BO{2026}
    \cycle{4}
    \competencesDuSocle{CA,RA}
    \objectifsApprentissage{C4D3ac,C4D1be}
    \motsCles{fonction,image,antécédent,Thalès,vitesse}
\enddatakeys
\begin{document}%

\medskip

On considère la fonction $f$ définie par $f(x) = - \dfrac{3}{19}x + 3$ pour tout nombre $x$ compris entre 0 et 19.

\medskip

\begin{enumerate}
\item Calculer l’image de 6 par la fonction $f$.
\item Déterminer l’antécédent de $0$ par la fonction $f$.
\end{enumerate}

\medskip

Un joueur de tennis effectue un service.

Voici une figure, qui n’est pas à l’échelle, représentant la trajectoire de la balle lors de ce service.

Le triangle ABC est rectangle en A.

Le triangle EDC est rectangle en E.

Les droites (AE) et (BD) se coupent au point C.

AB $= 3$ m ; AC $= 19$ m ; AE $= 12$ m

\begin{center}
\psset{unit=1cm}
\begin{pspicture}(14.5,3)
\pscircle(0.3,2){0.15}
\psline(0.3,1.85)(0.3,1)(0.5,0.3)
\psline(0.3,1)(0.2,0.3)
\psline(0.3,1.4)(0.6,1.6)(0.5,1.8)
\psline(0.3,1.4)(0,1.7)(0.2,2)
\pspolygon[linewidth=1.25pt](0.5,0.3)(13.8,0.3)(0.5,2.5)%ACB
\psline[linewidth=1.25pt](8.9,0.3)(8.9,1.13)%ED
\psframe(0.5,0.3)(0.75,0.55)\psframe(8.9,0.3)(9.15,0.55)
\psline[linewidth=1.25pt](0.2,2)(0.4,2.35)
\rput{60}(0.5,2.5){\psellipse(0,0)(0.25,0.1)}
\uput[d](0.5,0.3){A} \uput[ul](0.5,2.5){B} \uput[d](13.8,0.3){C} \uput[ur](8.9,1.13){D} \uput[d](8.9,0.3){E}
\end{pspicture}
\end{center}

Le joueur est positionné au point A.

On considère que la balle lancée en B effectue un trajet en ligne droite, qu’elle passe au-dessus du filet en D et qu’elle touche le terrain adverse en C.

La longueur DE représente la hauteur de la balle lorsque celle-ci passe au-dessus du filet.

Voici une représentation graphique de la fonction $f$ qui modélise la trajectoire de la balle lors de ce service. On rappelle que $f(x) = - \dfrac{3}{19}x + 3$, pour tout nombre $x$ compris entre 0 et 19.

\begin{center}
\psset{unit=0.6cm}
\begin{pspicture}(-1,-1)(20,5)
\psgrid[gridlabels=0pt,subgriddiv=1,gridwidth=0.2pt](0,0)(20,5)
\psaxes[linewidth=1.25pt,labelFontSize=\scriptstyle]{->}(0,0)(0,0)(20,5)
\psline[linewidth=1.25pt,linecolor=blue](0,3)(19,0)
\uput[ur](0,0){A}\uput[ur](19,0){C}\uput[ur](0,3){B}
\uput[d](19.5,0){$x$} \uput[l](0,4.75){$y$}
\end{pspicture}
\end{center}

Tout point de cette représentation graphique a pour abscisse $x$ et pour ordonnée $f(x)$ où $x$

est un nombre compris entre $0$ et $19$.

Dans le repère, le point A a pour coordonnées (0~;~0), le point B a pour coordonnées (0~;~3) et

le point C a pour coordonnées $(19~;~0)$.

\medskip

\begin{enumerate}[resume]
\item On considère le point d’abscisse 6 de la représentation graphique de la fonction $f$. Déterminer l’ordonnée de ce point.
\item On admet que la distance BC parcourue par la balle lors de ce service est d’environ $19,2$ m. Lors de ce service, la balle a mis $0,34$ seconde pour parcourir la distance BC.

Un commentateur affirme que la vitesse moyenne de la balle lors de ce service est de $208$ km/h. Cette affirmation est-elle vraie ?
\item Déterminer la hauteur de la balle, exprimée en mètre, lorsque celle-ci passe au-dessus du filet.
\end{enumerate}

\showthesource
\end{document}
